% ============================================================
% 分节契约
% purpose:      方法形式化描述,审稿人能据此复现
% allowed claims: 每个设计选择配一句"为什么不用 X"(design.md 差异表)
% forbidden claims: 无消融支撑的优越性断言(留给 Results)
% inputs:       03-method/design.md · 03-method/figures/
% 验证方式:     与 design.md 逐条对照;框架图来自 scibox-diagram
% ============================================================
\section{Method}\label{sec:method}

\subsection{Problem Formulation}
<形式化:记号、输入输出、优化目标。>

\begin{equation}
  <目标函数/核心公式,与 design.md 编号一致>
  \label{eq:objective}
\end{equation}

\subsection{<核心模块 A>}
<机制描述;设计选择的理由,含"为什么不用 X"。>

\subsection{<核心模块 B>}
<同上。>

\subsection{Overall Framework}
% 框架图由 scibox-diagram 产出(draw.io 可编辑源文件在 03-method/figures/),
% 导出 PDF 放 figures/,此处引用:
% \begin{figure}[t]
%   \centering
%   \includegraphics[width=\linewidth]{framework.pdf}
%   \caption{<方法总览。读图即可理解主流程,不依赖正文。>}
%   \label{fig:framework}
% \end{figure}
<端到端流程走读。>

% 伪代码(可选):
% \begin{algorithm}
%   <design.md 中的伪代码,algorithm/algorithmic 宏包按需加>
% \end{algorithm}
